\section{\textit{Feature Engineering}}

Tahap \textit{feature engineering} bertujuan untuk mengekstraksi dan merangkai parameter kunci agar datanya dapat diproses secara optimal oleh algoritma. Rekayasa fitur ini diterapkan pada variabel waktu (temporal) maupun teks kategori (kategorikal) yang belum terstruktur.
\begin{itemize}
    \item \textbf{Ekstraksi Waktu Kedatangan dan Keberangkatan}\\ Format penulisan pada kolom \texttt{waktu\_checkin} dan \texttt{waktu\_checkout} diuraikan (\textit{parsing}) menggunakan pustaka operasional waktu khusus (\texttt{pd.to\_datetime}) guna mengambil angka jam numeriknya saja. Hasil ekstraksi ini disimpan ke dalam fitur baru bernama \texttt{checkin\_hour} dan \texttt{checkout\_hour} dengan rentang nilai bulat 0 hingga 23.
    \item \textbf{Penanganan Nilai Kosong (\textit{Missing Values}) Numerik}\\ Kolom yang memuat angka durasi, seperti \texttt{lama\_menginap}, sering kali memiliki baris yang kosong atau tidak valid (\textit{null/NaN}). Data tersebut diubah ke dalam tipe numerik menjadi angka $0$.
    \item \textbf{Penanganan (\textit{Encoding}) Data Kategorikal}\\ Beberapa kolom memiliki bentuk teks klasifikasi, seperti \texttt{jenis\_sewa}, \texttt{metode\_pembayaran}, \texttt{status\_pasutri}, dan \texttt{status\_kewarganegaraan}. Untuk menangani sel yang kosong pada kolom-kolom tersebut, nilai yang hilang diganti secara seragam dengan label teks \texttt{'Unknown'}. Selanjutnya, seluruh kategori teks tersebut dikonversi menjadi urutan angka menggunakan metode standar \textit{Label Encoding} (misalnya Harian menjadi 0, Mingguan menjadi 1, dan Bulanan menjadi 2). Konversi angka dari label ini memastikan bahwa perhitungan matematis saat pelatihan klasifikasi dapat berjalan lancar tanpa terhambat oleh masalah tipe data \textit{string}. 
\end{itemize}